%% -*- coding:utf-8 -*-
%%%%%%%%%%%%%%%%%%%%%%%%%%%%%%%%%%%%%%%%%%%%%%%%%%%%%%%%%
%%   $RCSfile: errata.tex,v $
%%  $Revision: 1.1 $
%%      $Date: 2003/09/22 15:43:23 $
%%     Author: Stefan Mueller (DFKI)
%%    Purpose: 
%%   Language: LaTeX
%%%%%%%%%%%%%%%%%%%%%%%%%%%%%%%%%%%%%%%%%%%%%%%%%%%%%%%%%
%% $Log: errata.tex,v $
%% Revision 1.1  2003/09/22 15:43:23  stefan
%% started
%%
%%%%%%%%%%%%%%%%%%%%%%%%%%%%%%%%%%%%%%%%%%%%%%%%%%%%%%%%%


\documentclass[headings=small,bibliography=totoc,index=totoc,%
toc=left,
%oneside,% Stauffenburg Style
10pt,%    Stauffenburg Style 
cleardoublepage=empty % Stauffenburg Style, LaTeX prints heading
%,draft
]{scrbook}


\usepackage{etex}

%\usepackage{layout} % to visualize the layout of a document

\usepackage[latin1,utf8x]{inputenc}


\usepackage[T1]{fontenc}
\usepackage{times}


% für JPs cat
%\usepackage[T1]{tipa}
\usepackage{tipa}

\usepackage{amssymb}

\usepackage{multicol}

%% \setlength{\textwidth}{130mm}
%% %\setlength{\textheight}{187mm}
%% \setlength{\textheight}{192mm}
%% \setlength{\headsep}{5mm}
%% \setlength{\oddsidemargin}{17mm}
%% \setlength{\evensidemargin}{17mm}


%\raggedbottom
% Fußnoten unten, falls Seite nicht voll
%\usepackage[bottom,multiple]{footmisc} 


\deffootnote{5mm}{0mm}{\textsuperscript{\normalfont\thefootnotemark}}

\widowpenalty = 10000 % keine Hurenkinder

\displaywidowpenalty=10000\relax
\predisplaypenalty=-200\relax



\usepackage{jambox}
\usepackage{stauffenburg}

%\usepackage{array} % muß hier schon geladen werden wegen glossary
%\usepackage[hyper]{glossary}\makeglossary

\usepackage{german}
\usepackage[german]{varioref}
% do not stop and warn! This will be tested in the final version
%\vrefwarning


\usepackage[final]{epsfig}
\usepackage{graphicx}

\usepackage{tikz-qtree}
% has strange side effects
%\tikzset{every tree node/.style={align=left, anchor=north}}
\tikzset{every roof node/.append style={inner sep=0.1pt,text height=2ex,text depth=0.3ex}}

%% \usepackage{multind}
%% %\usepackage{myshowidx}

%% \makeindex{subind}
%% \makeindex{autind}
%% %\makeindex{wordind}

\newcommand{\mod}{\textsc{mod}\xspace}  % wegen beamer.cls nicht in abbrev.sty
\newcommand{\rel}{\textsc{rel}\xspace}  % wegen avm.sty in abbrev.sty

\usepackage[figuresright]{rotating}


\usepackage{article-ex,makros.2e,
de-date,de-commands,tree-dvips,pstricks,pst-node,
datefooter,lastpage,float,comment,my-theorems,soul}


\usepackage{my-gb4e-article}

% mit der Index-Version geht die Silbentrennung nicht
\renewcommand{\word}[1]{\emph{#1}}

\newcommand{\dom}{\textsc{dom}\xspace}

\renewcommand{\nom}{\textit{nom}}
\renewcommand{\gen}{\textit{gen}}
\renewcommand{\dat}{\textit{dat}}
\renewcommand{\acc}{\textit{acc}}


\begin{document}


\title{Zusätzliche Aufgaben für \emph{Head-Driven Phrase Structure Grammar: Eine Einführung}}

\maketitle

\setcounter{chapter}{21}


\section{Einleitung}

\begin{enumerate}
\item Schreiben Sie eine Phrasenstrukturgrammatik, mit der man u.\,a.\ die Sätze in (\mex{1})
      analysieren kann, die die Wortfolgen in (\mex{2}) aber nicht zulässt.
      \eal
      \ex[]{
      Ein Kind lacht.
      }
      \ex[]{
      Er gibt einem Kind den Ball.
      }
      \ex[]{
        Sie geben dem Kind kein Eis.
      }
      \ex[]{
      Er denkt an die Ferien.
      }
%       \ex[]{
%       Er wartet neben dem Bushäuschen auf ein Wunder.
%       }
      \zl
      \eal
      \ex[*]{
        Dem Kind lacht.
      }
      \ex[*]{
        Er gibt einer Kind den Ball.
      }
      \ex[*]{
        Er geben einem Kind den Ball.
      }
      \zl 
\end{enumerate}

\section{Der Formalismus}

\begin{enumerate}
\item Wie kann man Fahrräder mit Merkmalstrukturen modellieren?

\end{enumerate}


\section{Valenz und Grammatikregeln}

\begin{enumerate}
\item Geben Sie die Valenzlisten für folgende Wörter an:
      \eal
      \ex der
      \ex auf (in \emph{auf dem Baum})
      \ex zersägt
      \ex nachdenke
      \ex dürstet
      \zl
\end{enumerate}


\bibliography{biblio}

\end{document}


